\paragraph{3.2. Стандартные обозначения и часто встречающиеся функции}
\subparagraph*{3.2.1}
Покажите, что если функции $f(n)$ и $g(n)$ монотонно неубывающие, то таковыми же являются и функции $f(n) + g(n)$ и $f(g(n))$, а если вдобавок $f(n)$ и $g(n)$ неотрицательны, то монотонно неубывающей является и функция $f(n) \cdot g(n)$.

\subparagraph*{3.2.2}
Докажите уравнение (3.16).

\subparagraph*{3.2.3}
Докажите уравнение (3.19). Докажите также, что $n! = \omega(2^n)$ и $n! = o(n^n)$.

\subparagraph*{3.2.4 $\star$}
Является ли функция $\lceil \lg n \rceil!$ полиномиально ограниченной? А функция $\lceil \lg \lg n \rceil!$?

\subparagraph*{3.2.5 $\star$}
Какая из функций $\lg(\lg^* n)$ и $\lg^*(\lg n)$ является асимптотически большей?

\subparagraph*{3.2.6}
Покажите, что золотое сечение $\phi$ и сопряженное с ним $\hat{\phi}$ удовлетворяют уравнению $x^2 = x + 1$.

\subparagraph*{3.2.7}
Докажите по индукции, что $i$-е число Фибоначчи удовлетворяет уравнению
\[ F_i = \frac{\phi^i - \hat{\phi}^i}{\sqrt{5}}, \]
где $\phi$ — золотое сечение, а $\hat{\phi}$ — сопряженное с ним.

\subparagraph*{3.2.8}
Покажите, что из $k \ln k = \Theta(n)$ вытекает $k = \Theta(n / \ln n)$.
